考虑一块无吸收、局域各向异性的薄光学 \(q\)-板，入射单色傍轴光沿 \(+z\) 传播，角频率为 \(\omega\)。在忽略板内衍射的 Jones 近似下，板的两正交主轴具有相位延迟 \(\delta\)，局域主轴相对于 \(x\) 轴的取向为 \(\alpha(\phi)=q\phi+\alpha_0\)。仅研究避开中心取向缺陷的环形孔径，入射横向振幅具有归一化径向包络 \(R(r)\) 和角相位 \(e^{i\ell\phi}\)，其中 \(\ell\in\mathbb Z\)。定义
\[
\mathbf e_\sigma=\frac1{\sqrt2}\binom{1}{i\sigma},\qquad \sigma=\pm1,
\]
并约定其自旋角动量沿 \(z\) 的投影为 \(\sigma\hbar\)。设旋转矩阵 \(\mathsf R(\alpha)=\begin{psmallmatrix}\cos\alpha&-\sin\alpha\\\sin\alpha&\cos\alpha\end{psmallmatrix}\)，忽略整体传播相位。
(1) 从两主轴的相位延迟构造局域 Jones 矩阵 \(\mathsf J(\alpha,\delta)\)，证明它是幺正的；对入射圆偏振 \(\mathbf e_\sigma\) 求同旋向与反旋向的复振幅以及两者功率比例。
(2) 令入射态为 \(R(r)e^{i\ell\phi}\mathbf e_\sigma\)。推导反旋向分量获得的几何相位和轨道角动量本征值；由主轴是无向轴而不是有向箭头，证明允许的 \(q\) 取值条件，并讨论 \(q=0\) 与半整数 \(q\) 的区别。
(3) 对单光子及入射光功率 \(P\) 分别求出射光束的平均自旋、轨道和总角动量沿 \(z\) 的投影，进而求光板承受的平均光学力矩。解释 \(q=1\) 时即使自旋翻转也可以没有净力矩，并核对 \(q=1/2,\delta=\pi\) 的结果。
(4) 改用入射水平线偏振 \(\mathbf e_x\) 且角向相位均匀（\(\ell=0\)），在板后放置水平检偏器。推导任意延迟 \(\delta\) 的角向归一化强度 \(I_x(\phi)/I_0\)，在 \(\delta=\pi\) 时数出 \(q\ne0\) 的亮瓣数；计算全方位平均透过率，并解释无检偏器时为何不会出现这些强度亮瓣。
